\documentclass[exercice,a4]{thedocument}%
\usepackage{texmaths-tikz}%
\usepackage{pstricks}%
\usepackage{pst-bezier}%
\usepackage{pstricks-add}%
\usepackage{pst-eucl}%
\usepackage{pst-func}%
\usepackage{pst-tree}%
\usepackage{pst-circ}%
\usepackage{tikz}%
\usepackage{pgfplots}%
\usepackage{tkz-tab}%
\usepackage{tikzpagenodes}%
\usepackage{tabularray}%
\usepackage{tabularx}%
\usepackage{multirow}%
\usepackage{multicol}%
\usepackage{scratch3}%
\usepackage{enumitem}%
\usepackage[np]{numprint}%
\usepackage{pas-tableur}%
\usepackage{listings}%
\usepackage{graphicx}%
\usepackage{adjustbox}%
\usepackage{needspace}%
\usepackage{xcolor}%
\usepackage{eurosym}%
\usepackage{diagbox}%
\usepackage{colortbl}%
\usepackage{multido}%
\usepackage{icomma}%
\usepackage[french]{babel}%
\graphicspath{{../brevet/2026/}}%
\startdatakeys
    \source{26GENMATMEAG1}
    \BO{2026}
    \cycle{4}
    \competencesDuSocle{CA,RA,CO}
    \objectifsApprentissage{C4A1aa,C4A2ad,C4B2cb,C4D2aa,C4D3ac}
    \motsCles{réciproque de Pythagore,Thalès,angle,triangles semblables,aire}
\enddatakeys
\definecolor{gris}{gray}{0.5}%
\providecommand{\wideparen}[1]{\overset{\frown}{#1}}%
\begin{document}%
\providecommand{\degres}{^\circ}%
Sur la figure ci-dessous :

\begin{itemize}[left=5mm]
	\item les droites (BD) et (CE) sont sécantes en A;
	\item les droites (BC) et (DE) sont parallèles;
	\item $\mathrm{AB}=\np[cm]{6,4}$;\quad
	$\mathrm{AC}= \np[cm]{4,8}$;\quad
	$\mathrm{AD}= \np[cm]{4,8}$;\quad
	$\mathrm{BC}= \np[cm]{8}$.
\end{itemize}
\begin{center}
	\begin{tikzpicture}[baseline={(0,1.25)}]
		\draw[line width=1.2pt] (0,0)node[below left]{A}
		(6.4,0)node[right]{B}
		-- (0,4.8)node[above]{C} node[pos=0.5,above right]{\np[cm]{8}}
		-- (-4.8,0)node[left]{D}
		-- (0,-3.6)node[below]{E}--cycle
		(6.4,0)--(0,0) node[pos=0.5,below]{\np[cm]{6,4}}
		(-4.8,0)--(0,0) node[pos=0.5,below]{\np[cm]{4,8}}
		(0,4.8)--(0,0) node[pos=0.5,right]{\np[cm]{4,8}}
		(0,-3.6)--(0,0);
	\end{tikzpicture}
\end{center}

%%%% pstricks
%\begin{center}
%\psset{unit=1cm,linewidth=1.5pt}
%\def\xmin {-6}   \def\xmax {7} \def\ymin{-5}   \def\ymax{6}
%\begin{pspicture*}(\xmin,\ymin)(\xmax,\ymax)
%\pspolygon(-4.8,0)(0,4.8)(6.4,0)(0,-3.6)
%\psline(-4.8,0)(6.4,0) \psline(0,4.8)(0,-3.6)
%\uput[d](-2.4,0){$4,8$ cm} \uput[d](3.2,0){$6,4$ cm}
%\uput[r](0,2.4){$4,8$ cm} \uput[ur](3.2,2.4){$8$ cm}
%\uput[dl](0,0){A} \uput[r](6.4,0){B} \uput[u](0,4.8){C}
%\uput[l](-4.8,0){D}  \uput[d](0,-3.6){E}   
%\end{pspicture*}
%\end{center}

\begin{enumerate}
	\item Montrer que le triangle ABC est rectangle en A.

	\item Montrer que $\mathrm{DE}=\np[cm]{6}$ et $\mathrm{AE}= \np[cm]{3,6}$.

	\item Montrer que les angles $\widehat{\mathrm{ABC}}$ et $\widehat{\mathrm{ADE}}$ sont égaux.

	\textbf{Détailler la réponse en précisant la démarche.}

	\item Montrer que les triangles ABC et ADE sont semblables.

	\item Déterminer l'aire du quadrilatère BCDE.
\end{enumerate}


\needspace{10\baselineskip}
% saut de page si moins de 10 lignes disponibles
\medskip
\showthesource
\end{document}
